\par\SourceLeaf{144}
Eine subtile Wortkritik könnte vielleicht Gründe auffinden, die für die erste Hypothese sprechen würden, wobei ich mich aber um so weniger aufhalten will, da ich selbst diese Hypothese für zulässig nicht halten, kann, und zwar hauptsächlich aus dem Grunde, weil sonst P. mit seinen eignen Grundsätzen (oben S. [133]) in Widerspruch stehen würde. Ich glaube vielmehr, dass er, verleitet durch das \textsc{Kritter}’sche Gutachten, — oder, wie man auch sagen kann, und wie ich bald umständlich zeigen werde, durch seine unrichtige Auffassung dieses Gutachtens — die Zahl von 15 Pensionen wie eine unübersteigliche oder fast unübersteigliche Schranke betrachtet habe, unter deren Schutz er seinen Plan mit voller Sicherheit machen könne.

Einen Abglanz ähnlichen Vertrauens finde ich in dem Votum E.’s wieder, welches sich dadurch auszeichnet, dass es das einzige ist, welches der fatalistischen Zahl 15 erwähnt, und dessen wesentlicher Inhalt, so weit er hieher gehört, in folgendem besteht.

Die beiden F.’schen Fragepunkte 1 und 2. sagt E., müssten ihre Entscheidung allein durch den wirklichen Fonds und dessen Ertrag verglichen mit \textit{den erprobten} von H. P. sehr einleuchtend vorgelegten Grundsätzen erhalten. Es könne sein, dass die verlängerte Dauer der Waisenpensionen verursache, dass die höchste Zahl der Witwen von 15 überstiegen werden müsse, was doch als mit der Sicherheit der Kasse unverträglich schlechterdings nie geschehen dürfe. Er halte daher für rathsam, die Verlängerung der Waisenpensionen nur mit der ausdrücklichen Beschränkung zu bewilligen, \textit{so lange die Anzahl der Pensionen dadurch nicht über 15 gesteigert werde}. Diese Einschränkung scheine desto nothwendiger, weil die Berechnung der Verhältnisse der Kinder keine \textit{so gewisse} Erfahrung für sich habe, wie die beobachtete Proportion der Witwen. Ich verstehe dies so: E., dessen Auffassung des P.’schen Plans offenbar der \textsc{Sch.}’schen gerade entgegengesetzt ist, halte zwar nach der \textsc{Kritter}-P.’schen Theorie für gewiss, dass die Anzahl der Witwen nie über die berechnete Zahl (also 13) gehen könne; es sei aber nicht eben so gewiss, dass nicht mehr als 2 Waisenpensionen dazu kommen könnten, und für den Fall, dass diess doch geschehe, und die Gesammtzahl der Pensionen dadurch über 15 getrieben werden würde, müsste die Beschränkung der Waisenpensionen vorbehalten werden.

Aehnliche Besorgnisse, dass die Anzahl der Waisenpensionen zu geringe angeschlagen sein möchte, waren auch von einigen andern Votanten geäussert; indessen ist weder denselben in dem Finalbeschlusse Folge gegeben, noch haben sie sich durch die Erfahrung bisher bestätigt. In der That haben von Errichtung der Witwenkasse bis jetzt noch niemals mehr als zwei Waisenpensionen gleichzeitig bestanden, wobei ich jedoch nicht unbemerkt lassen will, dass während eines kurzen Theils des Jahrs 1776 die Anzahl auf 3 gestiegen sein würde, wenn die Erstreckung der Waisenpensionen bis zum vollendeten \(20^{\text{sten}}\) Jahre schon damals Statt gefunden hätte.

Ganz anders aber verhält es sich mit der \textit{Witwenzahl}. Die vermeinte Gewissheit, dass diese nicht über 13 steigen könne, ist durch die neuern Erfahrungen zerstört, indem sie schon einmal auf 22 gestiegen ist, und sogar der Durchschnittswerth während der letzten 8 Jahre etwas über 19 betragen hat. Schon lange vor 1794 war die Anzahl der Witwen (nemlich gleichfalls, ohne die Waisenpensionen mitzuzählen), einmal auf 14 gestiegen, und hatte während eines fünfjährigen Zeitraums (1774—1779) den Durchschnittswerth 13 behauptet. Diese Thatsache, die wohl dem Curator aber freilich nicht den votirenden Professoren bekannt sein konnte, hätte, aus dem allein zulässigen oben S. [130] Z. [4] angedeuteten Gesichtspunkte betrachtet, schon damals zum Beweise der Unrichtigkeit der P.-E.’schen Voraussetzung dienen können.

Da nun aber die Annahme der Zahl 13 für die höchste Witwenzahl auf dem \textsc{Kritter}’schen Gutachten beruhet, in welchem die Hälfte der Zahl der stehenden Ehen als maassgebend für das Maximum der Witwenzahl aufgestellt ist, so scheint der Schluss natürlich, dass Fehler in demselben sein möchten, und ich bin demnach an dem Punkt angelangt, wo ich dieses Gutachten selbst der Kritik unterwerfen muss.
